% ═══════════════════════════════════════════════════════════════════════════════
% Frozen Brane BC Complete Analysis - INCLUDE VERSION
% Auto-generated from frozen_brane_bc_complete.tex
% Labels prefixed with DL: for Derivation Library inclusion
% ═══════════════════════════════════════════════════════════════════════════════

\subsection*{Frozen on Brane via $\xi$-BC: Complete Analysis}
\label{DL:frozen-brane-bc}

\textbf{Status:} \tagDer{} \\
\textbf{Purpose:} Investigate whether boundary conditions in $\xi$ can create attraction
in the vortex--vortex interaction potential.

% --- Section 1: Model ---
\subsubsection*{1. Two Distinct ``Frozen'' Concepts}

\textbf{Radial Step Frozen:} $f(r) = \Theta(r-a)$ --- step function in radial direction.
Creates core energy, gives geometric coefficients ($4\pi/3$, $2\pi^2$), required for $m_p/m_e = 6\pi^5$.

\textbf{Brane-Frozen ($\xi$-BC):} Boundary conditions at $\xi = 0, \delta$:
Neumann ($\partial_\xi f = 0$), Robin ($\partial_\xi f = \kappa f$), or Dirichlet ($f = v$).

\textbf{These are independent concepts.}

% --- Section 2: BC from Action ---
\subsubsection*{2. Boundary Conditions from Action Variation}
\label{DL:bc-from-action}

Starting from bulk + boundary action $S = S_{\text{bulk}} + S_\partial$, variation gives:
\begin{equation}
\label{DL:eq:bc-variation}
\partial_\xi f\big|_{\text{boundary}} = V'(f)\big|_{\text{boundary}}
\end{equation}

\begin{center}
\begin{tabular}{@{}lll@{}}
\toprule
Boundary Action & Resulting BC & Status \\
\midrule
$V = 0$ & Neumann: $\partial_\xi f = 0$ & [Der] \\
$V = \frac{1}{2}\kappa f^2$ & Robin: $\partial_\xi f = \kappa f$ & [Der] \\
$V = \frac{1}{2}\lambda(f-v)^2$, $\lambda\to\infty$ & Dirichlet: $f = v$ & [Der] \\
\bottomrule
\end{tabular}
\end{center}

% --- Section 3: Mode Spectrum ---
\subsubsection*{3. Mode Spectrum Under Various BC}
\label{DL:mode-spectrum}

Eigenvalue problem: $-d^2\phi_m/d\xi^2 = \mu_m^2 \phi_m$ on $\xi \in [0, \delta]$.

\begin{center}
\begin{tabular}{@{}lccc@{}}
\toprule
BC Type & Zero Mode? & Lowest $\mu$ & Log term? \\
\midrule
Neumann--Neumann & YES & $\mu_0 = 0$ & Present \\
Dirichlet--Dirichlet & NO & $\mu_1 = \pi/\delta$ & Absent \\
Robin--Robin & NO & $\mu_1 > 0$ & Absent \\
Dirichlet (matching $v$) & YES$^*$ & const.\ $v$ & Present$^*$ \\
\bottomrule
\end{tabular}
\end{center}

% --- Section 4: Green's Function ---
\subsubsection*{4. Effective 2D Green's Function}

With $\xi$-weight $w(\xi)$:
\begin{equation}
G_{\text{eff}}(r) = \sum_m \frac{W_m^2}{N_m} G_m(r), \quad
W_m = \int_0^\delta w(\xi) \phi_m(\xi)\,d\xi
\end{equation}

Zero mode: $G_0(r) = -(1/2\pi)\ln(r/L)$. Massive: $G_m(r) = (1/2\pi)K_0(\mu_m r)$.

% --- Section 5: Sign Analysis ---
\subsubsection*{5. The No-Go Result}
\label{DL:bc-no-go}

\textbf{Theorem (BC Do Not Create Attraction) [Der]:}
For any BC (Neumann/Robin/Dirichlet) and weight (uniform/boundary):
\begin{equation}
\label{DL:eq:no-attraction}
V'_{\mathrm{lin}}(d) > 0 \quad \text{for all } d > 0
\end{equation}

\textbf{Proof:}
\begin{itemize}[nosep]
    \item Zero mode present: $d[\ln d]/dd = 1/d > 0$
    \item No zero mode: $d[K_0(\mu d)]/dd = -\mu K_1 < 0$, but enters $V$ with sign giving $V' > 0$
\end{itemize}
All terms positive $\Rightarrow$ $V' > 0$ always.

% --- Section 6: Comparison ---
\subsubsection*{6. Radial-Frozen vs.\ Brane-Frozen}
\label{DL:radial-vs-brane}

\begin{center}
\begin{tabular}{@{}lcc@{}}
\toprule
Aspect & Radial-Frozen $f(r)$ & Brane-Frozen ($\xi$-BC) \\
\midrule
Creates core energy? & YES & NO \\
$V_{\text{core}} \to +\infty$ at $d\to 0$? & YES & NO \\
Resolves no-go? & \textbf{YES} & \textbf{NO} \\
Essential for minimum? & YES & NO \\
\bottomrule
\end{tabular}
\end{center}

\textbf{Minimum mechanism:}
$V_{\text{core}} \to +\infty$ (radial-frozen, topology) +
$V_{\text{lin}} \to +\infty$ at large $d$ (log growth) $\Rightarrow$
minimum by continuity.

% --- Section 7: Verdict ---
\subsubsection*{7. Verdict}
\label{DL:bc-verdict}

\begin{itemize}[nosep]
    \item[\checkmark] BC from action variation: [Der]
    \item[\checkmark] Mode spectrum under BC: [Der]
    \item[\checkmark] $V'_{\mathrm{lin}} > 0$ for all BC: [Der]
    \item[$\times$] ``BC create attraction'': \textbf{FALSE}
\end{itemize}

\textbf{Conclusion:} Thick-brane BC provide scale $\delta$ and mode structure, but do NOT
create attraction. The minimum arises from topological core repulsion, not $\xi$-BC effects.
Chapter 2's radial-frozen approach is \textbf{validated and essential}.
